% ID: SEL-D0008-BD-E07
% ESTADO: validada. S009, 2026-11-07; tres pasadas del mismo autor y revisión matemática, didáctica y visual.
\textbf{a) Intervalo normal aproximado.}\par\nopagebreak[4]
La estimación es $\widehat p=0.15$. Con muestra aleatoria independiente y recuentos 150 y 850, usamos el intervalo aproximado para proporciones. Para confianza 95\,\%, $\alpha=0.05$ y $z_{0.975}\approx1.959963985$.
\[
\widehat p\pm z_{0.975}\sqrt{\frac{\widehat p(1-\widehat p)}n}
=0.15\pm1.959963985\sqrt{\frac{0.15(0.85)}{1000}}.
\]
\[
\boxed{IC_{\text{95\%}}\approx(0.127868891,\ 0.172131109).}
\]
En porcentajes, aproximadamente entre 12.79\,\% y 17.21\,\%. La confianza es cobertura del procedimiento para la proporción poblacional, bajo el modelo de muestreo del ejercicio.

\textbf{b) Tamaño mínimo y precisión del cuantil.}\par\nopagebreak[4]
Se mantiene la proporción piloto 0.15 y se exige $E<0.01$, un punto porcentual:
\[
z_{0.975}\sqrt{\frac{0.15(0.85)}n}<0.01
\iff n>\frac{z_{0.975}^2(0.1275)}{0.01^2}
\approx4897.859996.
\]
Con el cuantil preciso, el tamaño mínimo es \boxed{n=4898}. Se comprueba que $E(4897)>0.01$ y $E(4898)<0.01$.

\textbf{Si se trabaja con la tabla del examen:} usando $z=1.96$, el umbral es $n>4898.04$, por lo que resulta \boxed{n=4899\text{ con }z=1.96}. El redondeo del cuantil hace fallar por muy poco a 4898 en este segundo cálculo. 4899 es una elección conservadora que también cumple con el cuantil preciso.

Se muestran ambos procedimientos para no mezclar un tamaño mínimo calculado con un cuantil y una comprobación hecha con otro. Los criterios asociados no fijan un resultado numérico. El tamaño indicado es total y la desigualdad de error es estricta; no se sustituye la proporción dada por 0.5.
